\documentclass[10pt,letterpaper]{article}
\usepackage[utf8]{inputenc}
\usepackage[T1]{fontenc}
\usepackage[spanish]{babel}
\usepackage{mathptmx}
\usepackage[top=1.55cm, bottom=1.55cm, left=1.85cm, right=1.85cm]{geometry}
\usepackage{graphicx}
\usepackage{xcolor}
\usepackage{titlesec}
\usepackage{enumitem}
\usepackage{float}
\usepackage{csquotes}
\usepackage[backend=biber, style=ieee, sorting=none]{biblatex}
\addbibresource{referencias.bib}
\usepackage[hidelinks, plainpages=false, pdfpagelabels]{hyperref}

\titleformat{\section}{\normalfont\large\bfseries}{\thesection.}{0.5em}{}
\titleformat{\subsection}{\normalfont\normalsize\bfseries}{\thesubsection}{0.5em}{}
\titlespacing*{\section}{0pt}{6pt}{3pt}
\titlespacing*{\subsection}{0pt}{5pt}{2pt}

\setlength{\parindent}{0pt}
\setlength{\parskip}{3.5pt}

\begin{document}

% ==================== CARÁTULA ====================
\pagenumbering{gobble}
\begin{titlepage}
    \centering
    \includegraphics[width=3.2cm]{img/logo_uteq.png}\par
    \vspace{0.7cm}

    {\fontsize{16}{20}\selectfont\bfseries UNIVERSIDAD TÉCNICA ESTATAL DE QUEVEDO\par}
    \vspace{0.2cm}
    {\fontsize{13}{16}\selectfont\bfseries FACULTAD DE CIENCIAS DE LA COMPUTACIÓN\par}
    \vspace{0.15cm}
    {\fontsize{13}{16}\selectfont\bfseries CARRERA DE INGENIERÍA DE SOFTWARE\par}

    \vspace{1.2cm}
    {\fontsize{15}{19}\selectfont\bfseries APLICACIONES DISTRIBUIDAS\par}
    \vspace{0.25cm}
    {\fontsize{12}{15}\selectfont Séptimo Nivel -- Paralelo ``B''\par}

    \vspace{1.3cm}
    {\fontsize{14}{18}\selectfont\bfseries Investigación individual: arquitecturas distribuidas en la industria\par}
    \vspace{0.3cm}
    {\fontsize{11.5}{14}\selectfont\bfseries FORM-IC-01 -- Caso de estudio: Amazon (Microservicios)\par}

    \vspace{1.5cm}
    {\fontsize{12}{15}\selectfont\bfseries Estudiante:\par}
    \vspace{0.2cm}
    {\fontsize{12}{15}\selectfont MACIAS VÁSQUEZ CÉSAR ABRAHAN\par}

    \vspace{1.4cm}
    {\fontsize{12}{15}\selectfont\bfseries GitHub:\par}
    \vspace{0.25cm}
    \noindent\mbox{\href{https://github.com/CesarMacias007/FORM-IC-01-Investigaci-n-Individual-Arquitecturas-Distribuidas-en-la-Industria.git}{\textcolor{blue}{\ttfamily\scriptsize https://github.com/CesarMacias007/FORM-IC-01-Investigaci-n-Individual-Arquitecturas-Distribuidas-en-la-Industria.git}}}\par

    \vfill
    {\fontsize{12}{16}\selectfont\bfseries Docente: \normalfont Gleiston Cicerón Guerrero Ulloa, PhD\par}
    \vspace{0.25cm}
    {\fontsize{12}{16}\selectfont\bfseries Fecha: \normalfont 25 de noviembre de 2026\par}
    \vspace{0.25cm}
    {\fontsize{12}{16}\selectfont\bfseries Periodo Académico: \normalfont SPA 2026--2027\par}
\end{titlepage}

\pagenumbering{arabic}
% ==================== ANÁLISIS (1 PÁGINA EXACTA) ====================
\section*{Análisis del caso industrial: Microservicios en Amazon}

\subsection*{1. Descripción del caso}
Amazon opera su plataforma global de comercio electrónico (Amazon.com) y sus servicios de nube (AWS) bajo una arquitectura de \textbf{microservicios} evolucionada sobre infraestructura distribuida compartida. El estilo consiste en descomponer capacidades de negocio en servicios pequeños, autónomos y desplegables de forma independiente. La escala documentada en fuentes oficiales refleja volúmenes masivos: el servicio de almacenamiento Amazon S3 nació en 2006 con cerca de ocho servicios y creció hasta superar los \textbf{300 microservicios} especializados en políticas, clases de almacenamiento y gestión de datos \cite{vogels2023monoliths}. Durante el evento comercial Prime Day 2025, la infraestructura que respalda a Amazon.com registró picos de \textbf{151 millones de solicitudes por segundo} en Amazon DynamoDB y más de \textbf{1.7 billones ($1.7 \times 10^{12}$) de invocaciones diarias} en AWS Lambda, sosteniendo latencias de milisegundos de un solo dígito \cite{yun2025primeday}.

\subsection*{2. Motivo de la elección y arquitectura previa}
En sus primeros años, Amazon funcionaba sobre una aplicación monolítica escrita en C++ (conocida internamente como \emph{Obidos}) acoplada directamente a bases de datos centrales. Con el aumento del catálogo, el tráfico y el número de ingenieros, cualquier cambio menor obligaba a coordinar compilaciones y despliegues completos, mientras los accesos compartidos a base de datos generaban bloqueos severos. Guiada por su \emph{Distributed Computing Manifesto}, la empresa migró primero hacia una arquitectura orientada a servicios (SOA) y posteriormente hacia microservicios desacoplados \cite{vogels2023monoliths}. Esta transición buscaba alinear la estructura del código con equipos autónomos pequeños (regla de las ``dos pizzas''), eliminar dependencias rígidas en la capa de datos y permitir que cada componente escalara según su propia demanda \cite{velepucha2023survey}.

\subsection*{3. Beneficios obtenidos}
La adopción de microservicios aportó tres ventajas concretas a la operación de Amazon:
\begin{itemize}[leftmargin=*, itemsep=1pt, topsep=2pt]
    \item \textbf{Escalabilidad granular:} Permite ampliar únicamente los servicios sometidos a carga crítica sin replicar toda la plataforma, como se evidencia al absorber 151 millones de peticiones/s en DynamoDB durante Prime Day 2025 \cite{yun2025primeday}. Este comportamiento coincide con las mediciones experimentales de Blinowski et al., quienes confirman superioridad de los microservicios bajo cargas intensas y desbalanceadas \cite{blinowski2022monolithic}.
    \item \textbf{Evolución continua y aislamiento de fallos:} Un error en un servicio puntual acota su radio de impacto sin detener el proceso de compra completo, facilitando que S3 haya escalado a más de 300 microservicios a lo largo de casi dos décadas \cite{vogels2023monoliths}.
    \item \textbf{Velocidad de entrega:} Equipos independientes despliegan mejoras y eligen tecnologías específicas sin esperar ciclos monolíticos globales \cite{velepucha2023survey}.
\end{itemize}

\subsection*{4. Desafíos encontrados y soluciones aplicadas}
Operar cientos de servicios distribuidos trajo problemas no triviales. Primero, las llamadas síncronas encadenadas entre múltiples capas elevan el riesgo de sobrecarga en cascada cuando un servicio dependiente se vuelve lento. Amazon mitigó esto mediante técnicas de \textbf{descarte selectivo de carga} (\emph{load shedding}) para rechazar tráfico excedente antes de saturar recursos \cite{yanacek_loadshedding}, combinadas con \emph{timeouts} estrictos y reintentos con retroceso exponencial más fluctuación aleatoria (\emph{exponential backoff with jitter}) para evitar tormentas de reintentos simultáneos \cite{brooker_timeouts}. Segundo, la fragmentación introduce un ``impuesto de latencia'' por comunicación en red y mayor costo operativo, aspecto analizado por Blinowski et al. \cite{blinowski2022monolithic} y Velepucha y Flores \cite{velepucha2023survey}. Frente al dogmatismo de dividirlo todo, el propio CTO Werner Vogels recalcó en 2023 que un componente debe permanecer unificado cuando sus partes comparten perfil de escalado y un mismo equipo, como ocurrió con un servicio de monitoreo de Prime Video \cite{vogels2023monoliths}.

\subsection*{5. Aplicación al proyecto propio (PFC: TurnoJusto)}
Para el proyecto de fin de carrera \textbf{TurnoJusto: Turnos de Atención Distribuidos}, tomaría dos prácticas del caso Amazon: (a) delimitar responsabilidades por dominio (separando la gestión de colas/turnos en tiempo real del módulo de notificaciones o reportes) y (b) aplicar \emph{timeouts} acotados y reintentos con \emph{jitter} en las comunicaciones entre servicios para impedir bloqueos en cadena durante horas pico de atención. Por el contrario, \textbf{no} replicaría una fragmentación en decenas de microservicios ni el despliegue de múltiples bases de datos independientes desde el arranque. Dado que TurnoJusto atiende un volumen acotado y es desarrollado por un equipo pequeño, adoptar microservicios finos introduciría complejidad innecesaria en pruebas distribuidas, consistencia eventual y consumo de recursos \cite{blinowski2022monolithic, velepucha2023survey}; resulta más sensato partir de una arquitectura modular con dos o tres servicios bien justificados.

\clearpage

% ==================== EVIDENCIAS DE AUTORÍA ====================
\section*{Evidencias de autoría}

\subsection*{1. Captura de la fuente principal con fragmentos subrayados o anotados}
A continuación se presenta la captura de la fuente primaria oficial analizada (\emph{Monoliths are not dinosaurs}, publicado por Werner Vogels en \emph{All Things Distributed}), donde se destacan los fragmentos clave sobre la evolución de Amazon desde un monolito hacia microservicios y el crecimiento de Amazon S3 a más de 300 microservicios.

\begin{figure}[H]
    \centering
    \includegraphics[width=0.92\textwidth]{capturas/captura_fuente_principal.png}
    \caption{Captura anotada de la fuente principal oficial sobre arquitectura evolutiva y microservicios en Amazon \cite{vogels2023monoliths}.}
    \label{fig:captura_fuente}
\end{figure}

\subsection*{2. Sistema conocido de primera mano con el mismo estilo arquitectónico}
Un sistema que utilizo de manera cotidiana y en el que reconozco el mismo estilo de microservicios es \textbf{GitHub}. Como usuario y desarrollador, se observa que sus capacidades críticas operan como servicios independientes desacoplados: la gestión de repositorios Git, el servicio de autenticación y control de acceso, la plataforma de integración continua (\emph{GitHub Actions}), los \emph{webhooks} de eventos y la interfaz web de incidencias y \emph{pull requests}. En la práctica se comprueba el aislamiento de fallos propio de esta arquitectura: cuando la cola de ejecución de \emph{GitHub Actions} presenta demoras o degradación temporal, operaciones esenciales como realizar un \texttt{git push}, clonar código o revisar cambios en la web continúan funcionando con normalidad.

\subsection*{3. Declaración de uso de inteligencia artificial}
Se declara que se empleó un asistente de inteligencia artificial generativa como apoyo para estructurar el documento en \LaTeX, organizar las secciones solicitadas y verificar el formato de citas bibliográficas con \texttt{biblatex}. La selección del caso de Amazon, la lectura crítica de las fuentes oficiales, la comprobación manual de los DOI de las referencias académicas, el análisis aplicado al PFC \emph{TurnoJusto} y la revisión final del contenido fueron realizados y validados por el autor.

\clearpage

% ==================== REFERENCIAS ====================
\printbibliography[title={Referencias}]

\end{document}
